\documentclass[10.5pt]{article}

\usepackage[margin=0.85in]{geometry}
\usepackage{amsmath,amssymb,bm}
\usepackage{booktabs,array}
\usepackage{enumitem}
\usepackage{xcolor}
\usepackage{hyperref}
\usepackage{microtype}
\usepackage[skip=3pt,font=small]{caption}

\hypersetup{colorlinks=true,linkcolor=blue!55!black,citecolor=blue!55!black,urlcolor=blue!55!black}
\setlist{topsep=2pt,itemsep=1pt,parsep=0pt}

\newcommand{\ket}[1]{\left|#1\right\rangle}
\newcommand{\Var}{\operatorname{Var}}
\newcommand{\Cov}{\operatorname{Cov}}
\newcommand{\supp}{\operatorname{supp}}
\newcommand{\code}[1]{\texttt{#1}}
\newcommand{\sci}[2]{$#1\times10^{#2}$}

\title{Part II: Universal Measurable Estimators for ADAPT-VQE Gradient Selection\\[2pt]
\large Steps 0--4 and Phases 2--4 --- methods, results and reproduction}
\author{Part II working note}
\date{October 2026}

\begin{document}
\maketitle

\begin{abstract}
We implemented the Part II plan on the nine Part I benchmark cases and on states saved along exact ADAPT-VQE trajectories.

\begin{itemize}[leftmargin=1.2em]
  \item \textbf{Step 0.} Each measurement context reads out a maximal group of $2^n$ commuting Pauli products. Measurement circuits, the moments of every group from one transform, and exact shot-level sampling are validated against dense linear algebra and against Part I.
  \item \textbf{Step 1.} The II-0 baseline runs Part I's estimator as a shot-level experiment (pairwise elimination, $\gamma=0.9$, top-up allocation). Its median cost is 62--83\% below Part I's M3 as run.
  \item \textbf{Steps 2--3.} 60--76\% of $\mathcal B_0$ (mass completion) is measured by at least two contexts. With exact covariances, coefficient splitting (II-A) lowers the M3 cost by 38--80\%. A design fitted to Hartree--Fock covariances costs 66--69\% \emph{more} than II-0 on correlated states.
  \item \textbf{Step 4.} Learned from the shots, with estimated radii, II-A costs 35--57\% less than II-0 on the five fixed states run (H$_4$ CISD at two geometries, LiH, H$_2$O CISD at two geometries), with no wrong selections in 700 trials: 73--75\% of the online oracle ceiling on H$_4$ and LiH, and as much as the oracle design on H$_2$O. (The H$_4$ and LiH numbers replace the 36/43/39\% first reported, whose II-0 baselines predated the Step 4 fixes.)
  \item \textbf{Sign-aware elimination.} Part I's pairwise rule plugs in estimated gradient signs. On a correlated ADAPT state of stretched H$_4$ it selected the wrong generator in 5--24\% of trials for every configuration except II-A with oracle covariances, with heavy cost tails. A rule that uses pairwise contrasts only where signs are resolved selects correctly in every trial, at 21--31k shots. On fixed and trajectory states it costs $-2$ to $+16\%$ on H$_2$O, 6--38\% on H$_4$ and 5--80\% on LiH.
  \item \textbf{Contrast designs (II-E).} Designing the estimators of the decisive differences $s_lg_l-tg_i$ directly helps on the two hardest small-gap states (H$_4$ 1.0 ADAPT5 and LiH ADAPT5: $-32\%$ and $-36\%$ against arm-wise II-A), is neutral on H$_2$O ($-2$ to $+6\%$), and hurts on wide-gap H$_4$ states ($+11$ to $+39\%$ against arm-wise II-A with the same rule).
  \item \textbf{Non-oracle start.} Starting from an a priori bound instead of $\max_i|g_i|$ changes the cost by $-35$ to $+11\%$ on most states, but it is unreliable on the states with the smallest gradients. On H$_2$O eq CISD ($|g_{(1)}|=0.007$) the contrast design selects correctly in only 92\% of trials, and on LiH ADAPT5 in 99\% with a heavy tail: early decisions rest on near-deterministic contexts whose sample variance is zero.
  \item \textbf{Coverage.} With Part I's convention (Bonferroni over the pool, a fresh look every round), 4--29\% of trials have a round in which some interval misses the truth; selections stay correct because misses rarely involve the decisive pair. A union bound over rounds removes the misses at 1.8--2.6$\times$ the cost where the variance estimates are sound, but not where small samples underestimate them (78\% of trials on LiH ADAPT5 with the non-oracle start).
  \item \textbf{ADAPT trajectories (Phases 2 and 4).} Exact ADAPT reaches chemical accuracy in 10 (H$_4$ 1.0), 6 (LiH), 17 and 18 (H$_2$O) steps; stretched H$_4$ stalls at 3.2\,mHa. Every trajectory has exactly tied steps (spin partners), so selection must stop at a $\rho$-good leader. With measured selections ($\rho=0.1$, fully non-oracle), every H$_4$ trajectory reproduces exact ADAPT and 88 of 90 LiH trajectories do (two take one extra step); all reach the exact endpoint. II-A lowers the cumulative selection cost by 41--44\% against II-0 on H$_4$ and by 42\% on LiH, and II-E by 60\% on LiH. Tied steps dominate the cost.
\end{itemize}
All LiH and H$_2$O finite-shot runs were completed on the Trillium cluster and are included here.
\end{abstract}

\tableofcontents

%=====================================================================
\section{Problem and conventions}
\label{sec:problem}

At one ADAPT-VQE step, with state $\ket\psi$ and pool $\{G_i\}_{i=1}^N$, the task is to identify $i^\star=\arg\max_i|g_i|$, where
\[
  g_i=\langle C_i\rangle_\psi,\qquad C_i=[H,G_i]=\sum_\ell A_{i\ell}P_\ell ,
\]
using as few \emph{context-shots} as possible. One context-shot is one state preparation plus one measurement in one fully commuting (FC) context.

\paragraph{Cases.}
The nine Part I cases, rebuilt with the current Part I pipeline:
\begin{itemize}
  \item systems: H$_4$ square with side $1.0$ and $2.0$\,\AA; LiH at $R=3.0$\,\AA; H$_2$O at $R_{\mathrm{OH}}=0.9572$ and $1.75$\,\AA;
  \item states: HF and CISD (LiH HF only), plus the ADAPT trajectory states of Sec.~\ref{sec:trajectory};
  \item STO-3G, Jordan--Wigner, spin-orbital UCCSD pool ($N=26,92,140$ on $8,12,14$ qubits);
  \item Part I's validation gate, and Part I's deterministic first-fit FC partition of $\mathcal B_0=\bigcup_i\supp(C_i)$ into parent contexts $\mathcal M_\alpha$.
\end{itemize}
The rebuilt $|g_i|$ agree with Part I's committed values to $3\times10^{-10}$. Some H$_4$ rows differ in the sign of some $g_i$; this is an orbital-phase convention, and every quantity reported here is invariant under it.

\paragraph{Confidence and radii.}
A confidence radius is $r=z_\delta\epsilon$, with $\epsilon$ the one-standard-error target and $z_\delta=\Phi^{-1}(1-\delta/2N)$, $\delta=0.05$ (Bonferroni over the pool, as in Part I). The same $z_\delta$ is used for pairwise statements and in every round; Sec.~\ref{sec:coverage} measures what this convention delivers. The gap is $|g_{(1)}|-|g_{(2)}|$, and $\Delta_i=|g_{(1)}|-|g_i|$.

\paragraph{Oracle versus measured.}
Steps 2--3 are oracle: exact covariances drive the designs and exact gradients fix the elimination trajectory. Their numbers are ceilings. Steps 1 and 4 and Phases 2--4 sample measurement outcomes; exact quantities enter a non-oracle run only through its validation record.

%=====================================================================
\section{Measurement model (Step 0)}
\label{sec:measurement}

\paragraph{Measured groups.}
A context is measured by a Clifford $U_\alpha$ followed by readout of all $n$ qubits. One shot therefore gives the eigenvalues of all $2^n$ elements of the maximal abelian group $\mathcal G_\alpha=\{\pm U_\alpha^\dagger Z^{\bm z}U_\alpha\}$. A library Pauli $P_\ell$ is \emph{measured by} $\alpha$ if $P_\ell\in\mathcal G_\alpha$, equivalently if it commutes with all $n$ generators of $\mathcal G_\alpha$.

\paragraph{Completion.}
The parent members' generators are completed to $n$ by adding, one at a time, the first candidate that commutes with all current generators and lies outside their span. Candidate orders: \emph{canonical} (single-qubit $Z_q$, $X_q$, $Y_q$) or \emph{mass} (library Paulis by decreasing $\sum_i|A_{i\ell}|$); if a list runs out, a vector of the symplectic complement is used. A seeded \emph{random} order is a control, and ``span'' (no completion) a diagnostic lower bound on overlap.

\paragraph{Circuit synthesis.}
From the tableau $[X|Z]$ of the $n$ generators: row-reduce $X$ and apply $H$ on non-pivot qubits; row-reduce to $X=I$, which makes $Z$ symmetric; apply $S$ where $Z_{qq}=1$ and $CZ$ where $Z_{ab}=1$; apply $H$ on all qubits. Every Pauli of interest is then propagated through the gate list with exact phases, giving $U_\alpha P_\ell U_\alpha^\dagger=s_\ell Z^{\bm z_\ell}$.

\paragraph{Moments and sampling.}
With outcome distribution $p_\alpha(b)=|\langle b|U_\alpha|\psi\rangle|^2$, all group means come from one Walsh--Hadamard transform:
\begin{equation}
  E_\alpha[\bm z]=\sum_b p_\alpha(b)(-1)^{\bm z\cdot b},\qquad
  \mu_\ell=s_\ell E_\alpha[\bm z_\ell],\qquad
  \Cov(P_l,P_m)=s_ls_m\big(E_\alpha[\bm z_l\oplus\bm z_m]-E_\alpha[\bm z_l]E_\alpha[\bm z_m]\big).
  \label{eq:moments}
\end{equation}
The exact $p_\alpha$ gives the oracle moments; a histogram of $m$ sampled outcomes gives the empirical ones (factor $m/(m-1)$ on the covariance). Sampling $m$ shots is one multinomial draw over $2^n$ outcomes.

%=====================================================================
\section{Estimators and levels}
\label{sec:levels}

For a fixed context library, every linear unbiased estimator has the form
\begin{equation}
  \hat g_i=\sum_\alpha\bm x_{i\alpha}^{\!\top}\hat{\bm\mu}_\alpha,
  \qquad
  \sum_{\alpha:\,P_\ell\in\mathcal G_\alpha}x_{i\alpha\ell}=A_{i\ell}\quad\forall\ell .
  \label{eq:abc}
\end{equation}
With fragments $F_{i\alpha}=\sum_\ell x_{i\alpha\ell}P_\ell$, coefficient matrix $C$ and selection matrix $B$, Eq.~\eqref{eq:abc} is exactly $A=BC$. Every design, including every learned design at every refit, is checked for $\max|A-BC|\le10^{-9}\max|A|$ and the run fails hard otherwise.

The levels differ only in which coordinates $(\alpha,\ell)$ may be non-zero:
\begin{description}[leftmargin=3.4em,labelwidth=3em]
  \item[II-0] $\ell\in\supp(C_i)$ in its Part I parent context only: Part I's estimator.
  \item[II-A] $\ell\in\supp(C_i)$ in every context that measures it (coefficient splitting).
  \item[II-B] II-A plus Paulis with $A_{i\ell}=0$ measured by at least two contexts already carrying $C_i$ (zero-sum ghosts), capped at 200 per generator by a single-context control-variate score.
  \item[II-E] Direct contrast design: the target row is $s_lA_l-tA_i$ for a pair $(l,i)$, over the union of both arms' II-A coordinates. Pauli products whose coefficients cancel become zero-sum control variates (Sec.~\ref{sec:contrast}).
\end{description}
Coefficients are post-processing: any design satisfying Eq.~\eqref{eq:abc} can be applied to every shot already taken.

%=====================================================================
\section{Allocation, trajectories and cost metrics}
\label{sec:costs}

\paragraph{Allocation.}
With $n_\alpha$ shots per context, $\Var\hat g_i=\sum_\alpha q_{i\alpha}/n_\alpha$, where $q_{i\alpha}=\bm x_{i\alpha}^{\!\top}\Sigma_\alpha\bm x_{i\alpha}$. The allocation solves
\begin{equation}
  \min_{\bm n}\sum_\alpha n_\alpha\quad\text{s.t.}\quad\Var\hat g_i\le\epsilon^2\ \ \forall i\in\mathcal A,\qquad n_\alpha\ge c_\alpha .
  \label{eq:alloc}
\end{equation}
With $c=0$ this is Part I's solver. With committed shots $c\neq0$ (``top-up'') it is solved by multiplicative dual ascent, $n_\alpha=\max\!\big(c_\alpha,\sqrt{\sum_i\lambda_iq_{i\alpha}}\big)$, $\lambda_i\leftarrow\lambda_i(\Var_i/\epsilon^2)^{1/2}$, followed by a one-dimensional search for the smallest feasible scaling.

\paragraph{Cost metrics (Step 3).}
\begin{description}[leftmargin=5.8em,labelwidth=5.4em]
  \item[M1] One allocation over all $N$ arms at $r=\mathrm{gap}/2$.
  \item[M3 bound] Part I's exact-limit M3: at each breakpoint $\Delta_i$, arms with $\Delta_i\le2r$ are active and resolved to $r=\Delta/2$.
  \item[M3 actual] Part I's event-driven M3 on each arm's actual variance (primary metric).
\end{description}
\emph{Maxima} accounting (Part I) solves each breakpoint from scratch and charges each context its largest request; \emph{top-up} solves each breakpoint above the shots already committed. \emph{Static} designs are optimised for the full pool; \emph{redesign} (oracle II-D) re-optimises for the active set at every breakpoint.

%=====================================================================
\section{Joint design solver (Step 3)}
\label{sec:solver}

For an active set $\mathcal A$ the design problem $\min_{\bm x,\bm n}\sum_\alpha n_\alpha$ s.t.\ $\sum_\alpha q_{i\alpha}(\bm x)/n_\alpha\le\epsilon^2$ and Eq.~\eqref{eq:abc} is jointly convex (quadratic-over-linear) and scale-free. Alternating exact per-generator designs with the allocation (iterative coefficient splitting) is monotone but stalls: a context whose shots reach zero can never be refilled (on the H$_4$ 1.0 HF winner it ends at 2.590 against an optimum of 2.3766, in units of $\epsilon^{-2}$). Eliminating $\bm n$ gives the saddle problem
\[
  \min_{\bm x}\max_{\bm\lambda\ge0}\ \sum_\alpha h(Q_\alpha)-\epsilon^2\sum_i\lambda_i,
  \qquad Q_\alpha=\sum_i\lambda_iq_{i\alpha},
  \qquad h(Q)=n+Q/n,\ \ n=\max(c_\alpha,\sqrt Q),
\]
solved by an inner smoothed L-BFGS over all active arms' coefficients (smoothing $Q\to Q+\delta^2$, $\delta$ from $0.3$ to $10^{-7}$ of the typical context norm) and an outer multiplicative dual ascent. The reported cost is the exact allocation of the returned design. It matches an independent single-arm minimisation to $10^{-4}$ and SLSQP on small joint instances to $2\times10^{-4}$; on the large problems optimality is not certified.

%=====================================================================
\section{Online experiment (Steps 1 and 4)}
\label{sec:online}

\paragraph{Loop.}
Part I's successive elimination: the radius starts at $r_0$, shrinks as $r\leftarrow\gamma r$, and the run stops when one arm remains, when the leader is certified $\rho$-good (Sec.~\ref{sec:rules}), or when $r$ falls below a floor. Each round allocates shots for the active arms at $\epsilon=r/z$ (top-up), draws one multinomial histogram per context for the new shots, estimates $\hat g$, and eliminates. Every context that carries a fragment receives at least one shot (two in learned runs, one per fold). Trial $i$ uses the random stream \code{SeedSequence(seed).spawn(n)[i]}, so results do not depend on how trials are split across processes or machines.

\paragraph{Starting radius.}
\emph{Oracle start} (Part I): $r_0=\max_i|g_i|$, floor $10^{-3}r_0$. \emph{Bound start}: $r_0=\max_i\sum_\ell|A_{i\ell}|$, an a-priori bound on every $|g_i|$, floor $10^{-6}r_0$. With the bound start no exact quantity enters the run.

\paragraph{Learned designs (Step 4).}
\begin{itemize}
  \item \emph{Covariance model.} Per context and fold, $\widetilde\Sigma=\lambda\Sigma^{\rm prior}+(1-\lambda)\widehat\Sigma^{\rm emp}$, $\lambda=\nu/(\nu+m)$. Priors: the HF determinant, flat, none ($\nu=0$), or the oracle (diagnostic).
  \item \emph{Cross-fitting.} Every context's shots are split into two folds; the design fitted on one fold is applied to the other fold's means and the two estimates are averaged. Each half-estimate is unbiased whatever its design.
  \item \emph{Minimum data.} A context carries learned coefficients only when both folds hold at least 50 shots in it.
  \item \emph{Guard.} A fitted design is kept only if the other fold predicts a lower variance than the II-0 coefficients.
  \item \emph{Schedule.} An II-0 pilot round, then refits whenever the shots held have grown by 1.5.
  \item \emph{Radii and allocation} use only the held-out fold: $\Cov(\hat g_i,\hat g_j)=\tfrac14\sum_f\sum_\alpha\bm x^{f\top}_{i\alpha}\widetilde\Sigma^{\bar f}_\alpha\bm x^{f}_{j\alpha}/n^{\bar f}_\alpha$, and shots are planned with exactly this per-shot variance.
\end{itemize}
Three failures of the first implementation motivated the last three safeguards: fitting on one or two samples per fold (54\% correct on LiH), radii from a covariance including the fitting fold, and planning that disagreed with the radii (two H$_4$ 2.0 trials at $7.8\times10^8$ shots against a median of $1.3\times10^4$).

\paragraph{Small-sample radii.}
With the bound start, decisions can be taken when a context holds 2--5 shots per fold, where a sample covariance can be exactly zero. On a tied H$_4$ ADAPT state, II-0 with the bound start then stopped at 9{,}141 shots on an arm with 6\% of the best gradient. All configurations introduced in this note therefore replace the held-out covariance of a context with fewer than 50 shots by $k\,I$ over its $k$ Pauli products. This is a valid upper bound, since $\Sigma\preceq\operatorname{tr}(\Sigma)I\preceq kI$ when every Pauli has variance at most one, and it is used for radii and planning alike, so the planner buys such contexts out of the small-sample regime. The Step 4 rows of Table~\ref{tab:step4} predate this rule.

%---------------------------------------------------------------------
\subsection{Elimination rules and stopping}
\label{sec:rules}

Let $r_i=z\sqrt{\Var\hat g_i}$, $L_i=\max(0,|\hat g_i|-r_i)$, $U_i=|\hat g_i|+r_i$, and call the sign of $g_k$ \emph{resolved} when $|\hat g_k|>r_k$.
\begin{description}[leftmargin=5.4em,labelwidth=5em]
  \item[marginal] Arm $i$ is eliminated when $U_i<\max_jL_j$ (the work order's absolute-interval rule; valid whatever the signs).
  \item[pairwise] Part I's Eq.~(17): $i$ goes when some $j$ has $|\hat g_j|-|\hat g_i|>z\sqrt{\Var(s_j\hat g_j-s_i\hat g_i)}$, with $s=\operatorname{sign}\hat g$. The plug-in signs make it invalid while a sign is unresolved.
  \item[safe] $i$ goes when the marginal rule removes it, or when some $j$ with a resolved sign satisfies
  \[
    s_j\hat g_j-t\,\hat g_i>z\sqrt{\Var(s_j\hat g_j-t\,\hat g_i)}\qquad\text{for } t=s_i \text{ if $i$ is resolved, for both } t=\pm1 \text{ otherwise.}
  \]
  Given $s_j$, $|g_j|>|g_i|$ is equivalent to $s_jg_j>g_i$ and $s_jg_j>-g_i$, so the unresolved case needs no sign for $i$; on the event that all intervals hold, no step is wrong.
\end{description}
A test reproduces the failure of the plug-in rule: with $g=(0.020,-0.010)$, standard deviations $0.05$ and correlation $-0.999$, pairwise eliminates the best arm in over 20\% of draws ($z=2$), the safe rule in under 5\%.

\paragraph{$\rho$-good stopping.}
Exact BAI cannot separate an exact tie. With $\rho>0$ the run stops, returning the leader $l$, once $s_l\hat g_l-(1-\rho)t\hat g_i>z\sqrt{\Var(\cdot)}$ for $t=\pm1$ and every survivor $i$. Elimination stays exact, so on the good event the best arm survives and the returned arm has $|g_l|\ge(1-\rho)\max_i|g_i|$.

\paragraph{Anytime validity.}
Optionally, round $r$ uses $\delta_r=6\delta/(\pi^2r^2)$, which sums to $\delta$ over all rounds, so the family-wise statement covers the whole run instead of one round.

%---------------------------------------------------------------------
\subsection{Contrast objective (U\textsubscript{BAI}, II-E)}
\label{sec:contrast}

With the arm objective, each generator's estimator minimises its own variance and the shots are planned so that every active arm reaches $\epsilon$. At fixed shots the arms decouple, so this is also the minimiser of the work order's $\mathcal U_{\rm avg}=\sum_iw_i\Var\hat g_i$ for any weights.

The contrast objective targets what elimination actually tests. Once the empirical leader's sign is resolved, every other survivor $i$ receives a directly designed, cross-fitted estimator of $s_lg_l-tg_i$ ($t=s_i$ if resolved, both $t=\pm1$ otherwise), over the union of both arms' II-A coordinates. Its starting point and guard fallback is the difference of the two arm designs. Shots are planned so that each contrast's standard deviation reaches $2\epsilon$, the separation a pair of arms at radius $\epsilon$ would certify, and the leader's own reaches $\epsilon$. Designed contrasts take precedence over arm-implied ones in the safe rule. A test checks that contrast estimates are unbiased and that their held-out variances match the replicate spread (ratio 0.7--1.4). At II-0 the same construction changes only the allocation, which separates the allocation effect from the design effect.

%=====================================================================
\section{ADAPT trajectory states (Phase 2)}
\label{sec:trajectory}

Exact ADAPT-VQE is run from each HF case: $\ket{\psi_k}=e^{\theta_kG_{i_k}}\cdots e^{\theta_1G_{i_1}}\ket{\rm HF}$, the next generator being the one of largest $|g_i|$ (ties to the lowest index), with all parameters re-optimised by BFGS on an exact adjoint gradient. Everything is done in the $(N_\alpha,N_\beta)$ sector (36, 225 and 441 determinants), where $g_i=2\,\mathrm{Re}\langle H\psi|G_i\psi\rangle$. The run stops at $|E-E_{\rm FCI}|<1.6$\,mHa or when every $|g_i|<10^{-6}$.

Every state is saved and becomes a case \code{<geometry>\_ADAPT<k>} for all scripts. Its gradient problem is built by Part I's code from the \emph{same} Hamiltonian object as the trajectory: a fresh SCF can return orbitals of opposite sign to the cached HF problem (seen on H$_4$), which would flip some gradient signs relative to the saved states. Gradients rebuilt from the commutator expansions agree with the sector computation to $10^{-8}$ at every step.

\begin{table}[!ht]
\centering\small
\caption{Exact ADAPT trajectories. \emph{Ties}: steps whose two leading $|g_i|$ are equal (spin partners). \emph{Small gaps}: steps with relative gap below 5\%.}
\label{tab:trajectory}
\begin{tabular}{lrrll}
\toprule
Geometry & steps & final $|E-E_{\rm FCI}|$ & ties at step & small gaps at step \\
\midrule
H$_4$ side 1.0 & 10 & 1.09\,mHa & 2, 4, 8 & 1, 2, 4, 6, 8 \\
H$_4$ side 2.0 & 10 & 3.22\,mHa (stalls) & 1, 4, 8 & 1, 4, 6, 8 \\
LiH $R=3.0$ & 6 & 0.68\,mHa & 1, 4 & 1, 2, 4 \\
H$_2$O eq & 17 & 1.54\,mHa & 2, 4, 8 & 2, 3, 4, 6, 8, 9, 12, 14 \\
H$_2$O stretched & 18 & 1.43\,mHa & 2, 4 & 2, 4, 5, 7, 12, 13, 15, 16 \\
\bottomrule
\end{tabular}
\end{table}

On stretched H$_4$ every gradient vanishes by symmetry after 10 steps, 3.2\,mHa above FCI: a property of the spin-orbital pool, not of the measurement. Snapshots for fixed-state benchmarks were chosen mid-trajectory with distinct leaders (H$_4$ 1.0 ADAPT3 and ADAPT5, H$_4$ 2.0 ADAPT2 and ADAPT5, LiH ADAPT3 and ADAPT5, H$_2$O eq ADAPT11, H$_2$O stretched ADAPT8), plus tied states for $\rho$-good stopping.

%=====================================================================
\section{Phases 3 and 4: logged decisions and measured trajectories}
\label{sec:phases}

\paragraph{Phase 3 record.}
\code{scripts/phase3\_selection\_record.py} runs complete non-oracle decisions with a run log and writes the work order's required outputs: (i) $A$, labels and Hamiltonian; (ii) the library, multiplicities and auxiliary rule; (iii) every context with its Clifford circuit and images; (iv) the fold designs of every refit as sparse $C$ rows; (v) the $A=BC$ check at every refit (hard failure); (vi) outcome histograms per context and fold (the sufficient statistics); (vii) fold-model and prior covariance matrices of the most-used contexts; (viii)--(ix) per round, the shots added per context and fold, estimates, intervals, resolved signs, leader, eliminations with the arm and rule responsible, and designed contrasts; (x) context-shots, two-qubit gate counts and classical time; (xi) exact gradients and interval coverage in a separate \code{*\_validation.json}. Six logged decisions on LiH ADAPT3 (three per configuration, non-oracle start) were all correct, with no interval missing the truth in any round. II-0 (safe rule) took $1.09$--$1.27\times10^6$ shots in 47--49 rounds, and II-E took $4.8$--$6.8\times10^5$ shots in 52--53 rounds, with 13 refits and 6--10\,min of design time per decision.

\paragraph{Phase 4 driver.}
\code{scripts/phase4\_adapt\_selection.py} grows the ansatz with generators chosen by a Part II method from sampled outcomes on the current state; parameters are re-optimised exactly. Only selection shots are counted; energy estimation is outside Part II. Every method is fully non-oracle inside a selection (estimated radii, bound start, safe rule, small-sample radii) and stops at a $\rho$-good leader with $\rho=0.1$. The driver's own stopping rule (energy within 1.6\,mHa, or all exact $|g_i|<10^{-6}$, or 10 steps beyond exact ADAPT) is not measured.

%=====================================================================
\section{Results}
\label{sec:results}

\subsection{Steps 0--3 (unchanged)}

\begin{itemize}
  \item \emph{Step 0.} Circuits, group moments and the sampler are validated (sampler bias below $3.6\sigma$ over 2{,}000/400 replicates, variance and correlation errors within $5\sqrt{2/R}$ and $5/\sqrt R$, down to 86 shots in total).
  \item \emph{Step 1.} The online loop reproduces Part I's noiseless trial exactly (real-valued shots) and its Gaussian surrogate within 1.5 standard errors. II-0 (pairwise rule, $\gamma=0.9$, top-up, oracle radii) has median cost 285{,}250 / 39{,}754 / 160{,}459 / 22{,}136 / 299{,}479 on H$_4$ 1.0 HF / CISD, H$_4$ 2.0 HF / CISD and LiH, 62--83\% below Part I's M3 as run, with 5 wrong selections in 1{,}100 trials, all on the tied H$_4$ HF rows.
  \item \emph{Step 2.} Mass completion: multiplicity 1.95 (H$_4$), 2.95 (LiH), 3.10 (H$_2$O); 60--76\% of $\mathcal B_0$ splittable, carrying 62--89\% of the coefficient mass.
  \item \emph{Step 3.} On M3 actual, II-A lowers the cost by 38--54\% on H$_4$, 73--76\% on LiH and 70--80\% on H$_2$O (Table~\ref{tab:step3}). II-B adds 2--10 points on H$_4$ and nothing resolvable on LiH and H$_2$O. Redesign at every breakpoint adds 17--19 points on the depth-dominated H$_4$ HF rows and 15--16 points on stretched H$_2$O HF, where about 63\% of the static design's cost falls after the first elimination; it adds 1--3 points on LiH and the other H$_2$O states (11--25\% after the first elimination). It helps only with top-up accounting. HF-fitted designs cost 66--69\% more than II-0 on the CISD states.
\end{itemize}

\begin{table}[!ht]
\centering\small
\caption{Oracle ceilings (Step 3), mass completion. II-0: M3 on actual precision, context-shots. II-A, II-B: change against II-0 on the same metric. Redesign: oracle II-D at every elimination breakpoint, change against II-0 on the M3 bound, both with top-up accounting. The LiH and H$_2$O rows are the Oct 2 runs (\code{runs/step3\_oracle\_ceiling\_table5.csv}). The solver minimises the M3 bound, not M3 actual: the earlier solve of II-A on those five states reached the bound within 2 points but M3 actual within 11 points ($-70\%$, $-80\%$, $-75\%$, $-64\%$, $-68\%$), so II-A/II-B differences of that size are not significant.}
\label{tab:step3}
\setlength{\tabcolsep}{5pt}
\begin{tabular}{lrrrr}
\toprule
Case & II-0 & II-A & II-B & redesign (II-A / II-B) \\
\midrule
H$_4$ 1.0 HF & 569,020 & $-38\%$ & $-40\%$ & $-57\%$ / $-58\%$ \\
H$_4$ 1.0 CISD & 49,729 & $-54\%$ & $-56\%$ & $-53\%$ / $-56\%$ \\
H$_4$ 2.0 HF & 228,733 & $-40\%$ & $-50\%$ & $-59\%$ / $-60\%$ \\
H$_4$ 2.0 CISD & 31,553 & $-43\%$ & $-50\%$ & $-50\%$ / $-57\%$ \\
LiH HF & 333,272 & $-73\%$ & $-71\%$ & $-83\%$ / $-83\%$ \\
H$_2$O eq HF & \sci{1.32}{7} & $-80\%$ & $-83\%$ & $-83\%$ / $-83\%$ \\
H$_2$O eq CISD & \sci{2.44}{10} & $-76\%$ & $-78\%$ & $-80\%$ / $-81\%$ \\
H$_2$O str HF & \sci{1.46}{7} & $-70\%$ & $-68\%$ & $-88\%$ / $-88\%$ \\
H$_2$O str CISD & \sci{1.22}{8} & $-79\%$ & $-69\%$ & $-75\%$ / $-75\%$ \\
\bottomrule
\end{tabular}
\end{table}

The other Step 0--3 tables are in \code{part2\_estimator\_design/RESULTS.md} and in the previous version of this note.

\subsection{Step 4: designs learned from the shots}

\begin{table}[!ht]
\centering\small
\caption{Learned designs, pairwise rule, oracle start. Mean (change against II-0 with the same kind of radii) and correct selections, out of 200 on H$_4$ and 100 on LiH. The three baseline rows were rerun after the Step 4 fixes; the first version of this note compared the II-A rows with baselines run before them.}
\label{tab:step4}
\setlength{\tabcolsep}{3pt}
\begin{tabular}{llrrrrrr}
\toprule
 & & \multicolumn{2}{c}{H$_4$ 1.0 CISD} & \multicolumn{2}{c}{H$_4$ 2.0 CISD} & \multicolumn{2}{c}{LiH 3.0 HF} \\
\cmidrule(lr){3-4}\cmidrule(lr){5-6}\cmidrule(lr){7-8}
Design & radii & mean (change) & ok & mean (change) & ok & mean (change) & ok \\
\midrule
II-0 & oracle & 43,864 & 200 & 26,751 & 200 & 307,885 & 100 \\
II-A, oracle covariance & oracle & 24,317 ($-45\%$) & 200 & 11,643 ($-56\%$) & 200 & 146,538 ($-52\%$) & 100 \\
II-A, HF prior only & oracle & \sci{1.19}{9} & 200 & \sci{7.39}{8} & 194 & -- & -- \\
\midrule
II-0 & estimated & 41,463 & 200 & 24,699 & 200 & 304,990 & 100 \\
II-A, HF prior only ($\nu=\infty$) & estimated & 81,047 ($+95\%$) & 197 & 143,851 ($+482\%$) & 171 & -- & -- \\
II-A, HF shrunk, $\nu=1000$ & estimated & 32,099 ($-23\%$) & 200 & 31,696 ($+28\%$) & 195 & -- & -- \\
II-A, HF shrunk, $\nu=100$ & estimated & 28,307 ($-32\%$) & 200 & 15,293 ($-38\%$) & 200 & -- & -- \\
II-A, flat prior, $\nu=100$ & estimated & -- & -- & -- & -- & 193,042 ($-37\%$) & 100 \\
\textbf{II-A, data only ($\nu=0$)} & estimated & \textbf{27,129 ($-35\%$)} & \textbf{200} & \textbf{13,727 ($-44\%$)} & \textbf{200} & \textbf{184,511 ($-40\%$)} & \textbf{100} \\
II-A, data only, no guard & estimated & 26,749 ($-35\%$) & 200 & 13,764 ($-44\%$) & 200 & 184,511 ($-40\%$) & 100 \\
\bottomrule
\end{tabular}
\end{table}

\begin{itemize}
  \item \emph{Learned from the data, II-A keeps 73--75\% of its ceiling.} It costs 35\%, 44\% and 40\% less than II-0 run the same way; the online oracle ceiling is 45\%, 56\% and 52\%. All 500 trials selected the winner, and the data-only rows have no tail.
  \item \emph{The same holds on H$_2$O} (Table~\ref{tab:h2o}): learned II-A costs 57\% and 54\% less than II-0 on the equilibrium and stretched CISD states, as much as the design with exact covariances, with every selection correct.
  \item \emph{Estimated radii are nearly free at II-0} (within 6\% of oracle radii).
  \item \emph{The HF prior alone fails}, in the design and through near-zero variances in the radii and allocation. A prior helps only with little weight and never beat the data alone; the guard rarely acts once the minimum-data rule is in place.
\end{itemize}

\begin{table}[!ht]
\centering\small
\caption{H$_2$O CISD states, 100 trials each (two shards on Trillium), estimated radii unless marked. Mean context-shots (change against II-0, pairwise) and, where not 100, the percentage of correct selections in brackets.}
\label{tab:h2o}
\begin{tabular}{lrr}
\toprule
Configuration & H$_2$O eq CISD & H$_2$O stretched CISD \\
\midrule
II-0, pairwise & \sci{2.05}{10} & \sci{8.63}{7} \\
II-0, pairwise, oracle radii & \sci{2.07}{10} ($+1\%$) & \sci{9.19}{7} ($+6\%$) \\
II-A oracle covariance, pairwise & \sci{9.90}{9} ($-52\%$) & \sci{3.99}{7} ($-54\%$) \\
II-A HF $\nu=100$, pairwise & \sci{8.78}{9} ($-57\%$) & \sci{3.92}{7} ($-55\%$) \\
\textbf{II-A data, pairwise} & \sci{8.80}{9} ($-57\%$) & \sci{3.95}{7} ($-54\%$) \\
\midrule
II-0, safe & \sci{2.05}{10} ($+0\%$) & \sci{9.98}{7} ($+16\%$) \\
II-0, safe, bound start & \sci{2.02}{10} ($-2\%$) [99] & \sci{1.02}{8} ($+18\%$) \\
II-A data, contrast & \sci{8.65}{9} ($-58\%$) & \sci{3.92}{7} ($-55\%$) \\
II-A data, contrast, bound start & \sci{2.31}{9} ($-89\%$) [92] & \sci{2.55}{7} ($-70\%$) \\
\bottomrule
\end{tabular}
\end{table}

\subsection{Sign-aware rule, contrast objective and non-oracle start}

\begin{table}[!ht]
\centering\small
\caption{Fixed states, all configurations with estimated radii. Mean context-shots (change against II-0 with the pairwise rule and estimated radii, Table~\ref{tab:step4}); every row selected the winner in every trial. ``Bound'': non-oracle start. ``Anytime'': union bound over rounds. LiH: 100 trials. H$_2$O: Table~\ref{tab:h2o}.}
\label{tab:step4plus}
\setlength{\tabcolsep}{4pt}
\begin{tabular}{lrrr}
\toprule
Configuration & H$_4$ 1.0 CISD & H$_4$ 2.0 CISD & LiH 3.0 HF \\
\midrule
II-0, pairwise (reference) & 41,463 & 24,699 & 304,990 \\
II-0, safe & 47,037 ($+13\%$) & 26,126 ($+6\%$) & 549,094 ($+80\%$) \\
II-0, safe, bound start & 42,411 ($+2\%$) & 24,815 ($+0\%$) & 568,140 ($+86\%$) \\
II-0, contrast allocation & 61,080 ($+47\%$) & 35,031 ($+42\%$) & -- \\
\midrule
II-A data, pairwise & 27,129 ($-35\%$) & 13,727 ($-44\%$) & 184,511 ($-40\%$) \\
II-A data, safe & 32,814 ($-21\%$) & 18,962 ($-23\%$) & -- \\
II-A data, contrast & 45,569 ($+10\%$) & 21,114 ($-15\%$) & 238,484 ($-22\%$) \\
II-A data, contrast, bound start & 40,618 ($-2\%$) & 17,926 ($-27\%$) & 227,290 ($-25\%$) \\
II-A data, contrast, bound start, anytime & 74,255 ($+79\%$) & 38,420 ($+56\%$) & -- \\
\midrule
II-A oracle covariance, pairwise & 24,317 & 11,643 & 146,538 \\
II-A oracle covariance, contrast & 37,027 & 15,297 & 146,205 \\
\bottomrule
\end{tabular}
\end{table}

\begin{itemize}
  \item \emph{The bound start is cheap but not always safe.} Removing the last oracle input (the starting radius) changes the cost by $-35$ to $+11\%$ at matched rule and objective on most states. On H$_2$O eq CISD, the state with the smallest gradients ($|g_{(1)}|=0.007$, gap $0.002$), it selects correctly in 99\% (II-0) and 92\% (II-E) of trials, with intervals missing the truth in 76\% and 99\% of trials. The early, coarse rounds take decisions on contexts with a few tens of shots whose rare outcomes have not yet appeared, so their sample variance is zero; the $kI$ bound only covers contexts below 50 shots per fold.
  \item \emph{The safe rule costs $+6$ to $+13\%$ at II-0 and $+21$ to $+38\%$ at II-A on H$_4$, $0$ and $+16\%$ on H$_2$O CISD, and 80\% at II-0 on LiH.} The safe rows also carry the small-sample radius bound, so the two effects are not separated; on LiH many lightly sampled contexts sit under the conservative $kI$ bound for longer.
  \item \emph{The contrast objective does not pay on these fixed states.} On H$_4$ 1.0 CISD it is 39\% dearer than arm-wise II-A with the same rule, and even with oracle covariances it is 52\% dearer than the arm objective. Contrast planning charges the leader's own $\epsilon$ plus a $2\epsilon$ contrast for every survivor; when the gap is large, resolving each arm to $\epsilon$ is cheaper. On LiH the two objectives tie at the oracle level, and on H$_2$O CISD the learned contrast design is within 2\% of learned II-A.
\end{itemize}

\subsection{Interval coverage}
\label{sec:coverage}

The validation record counts, per trial, whether any interval used in any round (marginal for every active arm, and every designed contrast) missed the exact value.

\begin{table}[!ht]
\centering\small
\caption{Share of trials with at least one interval missing the truth in some round. Nominal family-wise level per round: $\delta=0.05$.}
\label{tab:coverage}
\begin{tabular}{lrrr}
\toprule
Configuration & H$_4$ 1.0 CISD & H$_4$ 2.0 CISD & LiH 3.0 HF \\
\midrule
II-0, pairwise, oracle radii & 0.22 & 0.25 & 0.14 \\
II-0, pairwise, estimated radii & 0.23 & 0.29 & 0.22 \\
II-A, oracle covariance & 0.12 & 0.19 & 0.08 \\
II-0, safe & 0.16 & 0.04 & -- \\
II-A data, safe & 0.09 & 0.04 & -- \\
II-A data, contrast, bound start & 0.12 & 0.14 & -- \\
II-A data, contrast, bound start, anytime & 0.00 & 0.00 & -- \\
\bottomrule
\end{tabular}
\end{table}

Part I's convention does not deliver family-wise coverage over a whole run: a fresh look every round with a per-round Bonferroni bound misses in 4--29\% of trials (19--21\% for II-0 on H$_2$O CISD). The selections are nevertheless correct, because misses rarely involve the decisive pair. The anytime bound removes every miss on H$_4$ and LiH ADAPT3, at 1.8--2.6$\times$ the cost (Tables~\ref{tab:step4plus}, \ref{tab:snapshots} and~\ref{tab:snapshots2}). It does not on LiH ADAPT5 with the bound start (78\% of trials still miss): there the misses come from underestimated variances, not from repeated looks. On the trajectory state H$_4$ 2.0 ADAPT2 the learned pairwise configurations miss in 63--78\% of trials and the learned safe configurations in 1--2\%.

\subsection{Phase 2: correlated ADAPT states}

\begin{table}[!ht]
\centering\footnotesize
\caption{ADAPT snapshots of H$_4$, 100 trials each, estimated radii unless marked. Mean context-shots and, where not 100, the percentage of correct selections in brackets. $^\dagger$Tail-dominated mean; medians: H$_4$ 1.0 ADAPT5 II-0 oracle radii 558{,}631; H$_4$ 2.0 ADAPT2 II-0 26{,}920, II-0 oracle radii \sci{9.05}{7}, II-A data 17{,}770, HF prior 5{,}610, HF $\nu=100$ 7{,}807.}
\label{tab:snapshots}
\setlength{\tabcolsep}{3pt}
\begin{tabular}{lrrrr}
\toprule
 & \multicolumn{2}{c}{H$_4$ side 1.0} & \multicolumn{2}{c}{H$_4$ side 2.0} \\
\cmidrule(lr){2-3}\cmidrule(lr){4-5}
Configuration & ADAPT3 (gap 33\%) & ADAPT5 (gap 15\%) & ADAPT2 (gap 11\%) & ADAPT5 (gap 56\%) \\
\midrule
II-0, pairwise & 26,504 & 411,200 & 789,821$^\dagger$ [95] & 58,402 \\
II-0, pairwise, oracle radii & 26,455 & \sci{2.92}{8}$^\dagger$ [99] & \sci{6.70}{7}$^\dagger$ [92] & 53,636 \\
II-A oracle covariance, pairwise & 14,240 & 239,815 & 16,622 & 26,895 \\
II-A data, pairwise & 18,161 & 288,503 [99] & 246,207$^\dagger$ [95] & 29,749 \\
II-A HF prior only, pairwise & 16,954 & 455,298 [99] & 39,181$^\dagger$ [76] & 69,934 \\
II-A HF $\nu=100$, pairwise & 17,751 & 246,442 & 80,193$^\dagger$ [85] & 32,520 \\
\midrule
II-0, safe & 29,883 & 406,671 & 28,031 & 55,818 \\
II-0, safe, bound start & 28,233 & 399,827 & 30,726 & 53,882 \\
II-A data, contrast & 27,787 & \textbf{197,352} & 23,357 & 34,881 \\
II-A data, contrast, bound start & 24,513 & 218,787 [99] & 24,023 & 30,269 \\
II-A data, contrast, bound start, anytime & 51,251 & 568,295 & 62,195 & 62,772 \\
II-A oracle covariance, contrast & 17,895 & 193,558 & 21,049 & 28,543 \\
\bottomrule
\end{tabular}
\end{table}

\begin{itemize}
  \item \emph{The plug-in sign rule fails on correlated trajectory states.} On H$_4$ 2.0 ADAPT2 every pairwise configuration except II-A with oracle covariances made wrong selections (II-0 95\%, II-0 with \emph{oracle} radii 92\% and a $6.7\times10^7$ mean, II-A data 95\%, HF prior 76\%). With the safe rule every configuration selected correctly in every trial, at 21--31k shots (62k with the anytime bound). Near-zero gradients with unresolved signs and strong correlations are exactly the case the plug-in signs get wrong.
  \item \emph{On hard states the contrast objective is the best non-oracle method.} On H$_4$ 1.0 ADAPT5 (gap 15\%) it costs 197{,}352, 32\% below arm-wise II-A data and 52\% below II-0, and matches the oracle contrast design.
  \item \emph{Splitting still pays on trajectory states.} With the safe rule and the bound start, II-A contrast costs 13--45\% less than II-0.
  \item \emph{The HF prior is unreliable along the trajectory}: the cheapest non-oracle row on ADAPT3, but 76\% correct on ADAPT2.
\end{itemize}

\paragraph{Tied states.}
On H$_4$ 1.0 ADAPT4 (exact tie) exact BAI ran to $1.66\times10^9$ shots in a single trial. With $\rho=0.1$, II-0 safe (bound start) stops at a mean of 493{,}275 and II-A contrast at 236{,}115 shots; on H$_4$ 2.0 ADAPT1, 24{,}633 and 17{,}361. Every selection was $\rho$-good: 36--52\% picked the lower-indexed generator of the tie, 44--54\% its tie partner, and 4--10\% a third generator within 3--5\% of the best.

\begin{table}[!ht]
\centering\footnotesize
\caption{ADAPT snapshots of LiH and H$_2$O, 100 trials each, estimated radii unless marked. Mean context-shots (change against II-0, pairwise) and, where not 100, the percentage of correct selections in brackets. ``--'': not run. Relative gaps: LiH ADAPT3 75\%, ADAPT5 15\%, H$_2$O eq ADAPT11 11\%, stretched ADAPT8 15\%.}
\label{tab:snapshots2}
\setlength{\tabcolsep}{3pt}
\begin{tabular}{lrrrr}
\toprule
Configuration & LiH ADAPT3 & LiH ADAPT5 & H$_2$O eq ADAPT11 & H$_2$O str ADAPT8 \\
\midrule
II-0, pairwise & 709{,}261 & \sci{6.60}{7} & \sci{1.93}{8} & \sci{1.84}{8} \\
II-0, pairwise, oracle radii & 707{,}341 ($-0\%$) & \sci{6.60}{7} ($+0\%$) & -- & -- \\
II-A oracle covariance, pairwise & 375{,}474 ($-47\%$) & \sci{3.61}{7} ($-45\%$) & \sci{7.43}{7} ($-62\%$) & \sci{9.03}{7} ($-51\%$) \\
II-A data, pairwise & 430{,}494 ($-39\%$) & \sci{4.34}{7} ($-34\%$) & \sci{8.17}{7} ($-58\%$) & \sci{8.57}{7} ($-53\%$) \\
II-A HF prior only, pairwise & 437{,}864 ($-38\%$) & \sci{2.36}{7} ($-64\%$) [29] & -- & -- \\
II-A HF $\nu=100$, pairwise & 430{,}152 ($-39\%$) & \sci{4.14}{7} ($-37\%$) & \sci{7.56}{7} ($-61\%$) & \sci{8.61}{7} ($-53\%$) \\
\midrule
II-0, safe & \sci{1.14}{6} ($+61\%$) & \sci{6.92}{7} ($+5\%$) & \sci{2.01}{8} ($+4\%$) & \sci{1.81}{8} ($-2\%$) \\
II-0, safe, bound start & \sci{1.21}{6} ($+71\%$) & \sci{6.64}{7} ($+1\%$) & \sci{1.92}{8} ($-1\%$) & \sci{1.85}{8} ($+1\%$) \\
II-A data, contrast & 600{,}144 ($-15\%$) & \sci{2.77}{7} ($-58\%$) & \sci{8.24}{7} ($-57\%$) & \sci{9.10}{7} ($-50\%$) \\
II-A data, contrast, bound start & 530{,}004 ($-25\%$) & \sci{1.16}{8} ($+76\%$) [99] & \sci{6.66}{7} ($-66\%$) & \sci{7.88}{7} ($-57\%$) \\
II-A data, contrast, bound start, anytime & 910{,}289 ($+28\%$) & \sci{1.93}{8} ($+192\%$) & -- & -- \\
II-A oracle covariance, contrast & 370{,}494 ($-48\%$) & \sci{2.91}{7} ($-56\%$) & \sci{7.81}{7} ($-60\%$) & \sci{9.44}{7} ($-49\%$) \\
\bottomrule
\end{tabular}
\end{table}

\paragraph{LiH and H$_2$O snapshots (Table~\ref{tab:snapshots2}).}
\begin{itemize}
  \item Learned splitting pays on every state: II-A data costs 34--39\% less than II-0 on LiH and 53--58\% less on H$_2$O, with every selection correct.
  \item The contrast design is the best non-oracle method on LiH ADAPT5 (gap 15\%), 36\% below II-A data and level with its oracle form. On H$_2$O it is within $+1$ to $+6\%$ of II-A, and on the wide-gap LiH ADAPT3 it is 39\% dearer.
  \item The HF prior alone selects correctly in only 29\% of trials on LiH ADAPT5. Shrunk to the data ($\nu=100$), it matches the data alone within 8\% on all four states.
  \item The bound start makes the contrast design 13--19\% cheaper on H$_2$O, but on LiH ADAPT5 it costs a 1\% error rate and a heavy tail (median $5.4\times10^7$, mean $1.2\times10^8$).
  \item On the tied LiH ADAPT4, $\rho=0.1$ stopping costs $4.1\times10^7$ (II-0) and $1.6\times10^7$ shots (II-E, $-62\%$). 100\% and 97\% of the selections are $\rho$-good, so II-E's 3\% exceeds $\delta$, for the same small-sample reason.
\end{itemize}

\subsection{Phase 4: measured ADAPT trajectories}

\begin{table}[!ht]
\centering\small
\caption{Measured ADAPT trajectories (48 per method on H$_4$, 30 on LiH), $\rho=0.1$, fully non-oracle selections. Cumulative selection context-shots over the whole trajectory. Every trajectory ended at the exact-ADAPT energy; all H$_4$ and 88 of 90 LiH trajectories took the same number of steps. ``Best'': selections whose $|g|$ equals the maximum (ties included).}
\label{tab:phase4}
\begin{tabular}{llrrrr}
\toprule
Geometry & method & median & P90 & best & $\rho$-good \\
\midrule
H$_4$ side 1.0 & II-0, safe & \sci{1.88}{6} & \sci{2.34}{6} & 93\% & 100\% \\
(10 steps, 1.09\,mHa) & II-A data, safe & \sci{1.11}{6} ($-41\%$) & \sci{1.44}{6} & 92\% & 100\% \\
 & II-A data, contrast & \sci{1.09}{6} ($-42\%$) & \sci{1.61}{6} & 90\% & 100\% \\
\midrule
H$_4$ side 2.0 & II-0, safe & \sci{1.59}{7} & \sci{2.29}{7} & 94\% & 100\% \\
(10 steps, stalls at 3.22\,mHa) & II-A data, safe & \sci{9.39}{6} ($-41\%$) & \sci{1.56}{7} & 93\% & 100\% \\
 & II-A data, contrast & \sci{8.85}{6} ($-44\%$) & \sci{1.56}{7} & 94\% & 100\% \\
\midrule
LiH & II-0, safe & \sci{8.57}{7} & \sci{1.16}{8} & 94\% & 100\% \\
(6 steps, 0.68\,mHa) & II-A data, safe & \sci{4.94}{7} ($-42\%$) & \sci{7.66}{7} & 94\% & 99.4\% \\
 & II-A data, contrast & \sci{3.41}{7} ($-60\%$) & \sci{5.52}{7} & 96\% & 99.4\% \\
\bottomrule
\end{tabular}
\end{table}

\begin{itemize}
  \item Measured selections reproduce exact ADAPT: same length and final energy in all 288 H$_4$ trajectories; the mean shortfall of the selected $|g|$ is 0.13--0.20\%. On LiH one selection per learned method (of 181) was not $\rho$-good (shortfall 13--15\%); each such trajectory took one extra step and still reached the same energy.
  \item II-A lowers the cumulative selection cost by 41--44\% against II-0 at the same rule and stopping, and II-E by 42--60\%, most on LiH.
  \item The cost is dominated by tied and small-gap steps: on H$_4$ 2.0 the tied step 8 ($\max|g_i|=0.0094$) takes $7$--$14\times10^6$ of the total, and on H$_4$ 1.0 the tied steps 4 and 8 take about half. Without $\rho$-good stopping these steps would not terminate.
\end{itemize}

\subsection{Cluster runs}

The LiH snapshots, the tied ADAPT4 state, LiH Phase 4 and Phase 3, and all H$_2$O runs ran on Trillium: two LiH node-jobs (2.2 and 5.1\,h) and eight H$_2$O node-jobs (11--17.5\,h, 60 workers each, memory-bound), about 110 node-hours in total. None needed a requeue. On H$_2$O one II-A trial costs about 6\,min of CPU time, II-0 with the bound start 15\,min, and a learned contrast trial 1.4--5.5\,h of design time.

%=====================================================================
\section{Conclusions and next steps}

\begin{enumerate}
  \item The shot-level II-0 baseline is 62--83\% cheaper than Part I's M3 as run, and coefficient splitting has a large oracle ceiling (38--80\% on M3 actual).
  \item Learned from the shots, II-A costs 35--57\% less than II-0 without oracle input (H$_4$ CISD, LiH, H$_2$O CISD) and 30--58\% less on ADAPT trajectory states, provided it uses cross-fitting, held-out radii and allocation, and a minimum number of shots per fold. That is 73--75\% of the oracle ceiling on H$_4$ and LiH and all of it on H$_2$O. Data-driven covariances should be the default; the HF prior should never be used alone (29\% correct on LiH ADAPT5).
  \item Part I's plug-in sign rule is not safe on correlated ADAPT states (5--24\% wrong selections on H$_4$ 2.0 ADAPT2). The sign-aware rule fixes this at $-2$ to $+16\%$ extra cost on H$_2$O and 6--38\% on H$_4$, but up to 80\% on LiH; separating the rule from the small-sample radius bound on LiH is the next measurement.
  \item The starting-radius oracle can be removed on most states at no cost, but not yet safely on the hardest ones (92\% correct on H$_2$O eq CISD): variance estimates from a few tens of shots of near-deterministic contexts need a better upper bound (for example a Bernoulli upper confidence bound) before the non-oracle start can be the default.
  \item Direct contrast design (II-E) is the method of choice on hard, small-gap states ($-32\%$ and $-36\%$ against II-A; $-60\%$ against II-0 over LiH trajectories), neutral on H$_2$O, and should not be used when the gap is large. A natural rule is to switch from arm to contrast planning once few arms survive.
  \item Per-round Bonferroni does not give run-wide coverage. Report it as such, or pay 1.8--2.6$\times$ for the anytime bound, which is only as good as the variance estimates it is applied to.
  \item ADAPT trajectories contain exact ties at every geometry; $\rho$-good stopping is required, and with it measured selection reproduces exact ADAPT. Ties dominate the cumulative cost, so the next gains are in how quickly a tie is certified $\rho$-good.
\end{enumerate}

\paragraph{Next.}
Make small-sample variance bounds valid for near-deterministic contexts; run measured trajectories on H$_2$O; decompose the LiH cost of the safe rule; switch between arm and contrast objectives adaptively; replace per-generator refits by the joint solver; price circuits (8--32 CZ per context) and classical design time (now logged per run).

%=====================================================================
\section{Reproduction}
\label{sec:reproduction}

\paragraph{Code.}
\code{part2\_estimator\_design/} on branch \code{part2}; it imports Part I's \code{src/} directly. New modules: \code{rules} (elimination and stopping), \code{trajectory} (ADAPT states), \code{runlog} (Phase 3 record); \code{learning} gained the contrast objective, the bound start, the small-sample radius bound and anytime validity; \code{parallel} gained checkpoints.

\paragraph{Environment.}
Local: WSL Ubuntu, Python 3.12.3, NumPy 2.5.2, SciPy 1.18.1, PySCF 2.14.0, OpenFermion 1.8.1 (\code{requirements.txt}). Trillium: Python 3.12.4 with NumPy 2.4.2, SciPy 1.17.1, PySCF 2.14.0, OpenFermion 1.8.1 from the Alliance wheelhouse. Per-trial seeding makes results independent of process layout; they are statistically, not bitwise, comparable across library versions.

\paragraph{Commands}
(from \code{part2\_estimator\_design/}; \code{OMP\_NUM\_THREADS=1} with \code{--workers}):
\begin{verbatim}
python -m pytest tests -q                               # 67 tests
python scripts/phase2_adapt_states.py --cases H4_square_eq_side1p0_HF \
       H4_square_stretch_side2p0_HF LiH_R3p0_HF H2O_eq_HF H2O_stretch_HF
python scripts/step4_learned_designs.py --cases H4_square_eq_side1p0_CISD \
       --configs "II-0 safe" "II-A data, contrast" ... --trials 200 --workers 12
python scripts/step4_learned_designs.py --cases H4_square_eq_side1p0_ADAPT5 \
       --configs ... --trials 100 --workers 12
python scripts/phase3_selection_record.py --cases LiH_R3p0_ADAPT3 --trials 3
python scripts/phase4_adapt_selection.py --cases H4_square_eq_side1p0_HF \
       --trajectories 48 --workers 12
python scripts/merge_trials.py --step step4 --cases <cases>
python scripts/aggregate_results.py
\end{verbatim}
The Step 0--4 commands of the previous version are unchanged. Configuration names are the keys of \code{CONFIGS} in \code{scripts/step4\_learned\_designs.py}.

\paragraph{Cluster.}
\code{cluster/trillium\_run.sbatch} runs one whole node per job file in \code{cluster/jobs/}; the lines of a job file start together. Every finished trial and trajectory is checkpointed; with \code{--resume} a rerun skips finished configurations and trials, and the job requeues itself 15\,min before the 24\,h limit (at most twice). A resumed run reproduces an uninterrupted one exactly (tested). H$_2$O jobs are memory-bound at 60 workers per node; LiH jobs use about 190 workers.

\paragraph{Outputs.}
Each case writes \code{runs/<case>/<case>\_<step>.csv} (rebuilt from all per-trial files, so runs of different configurations accumulate), \code{runs/<case>/trials/}, \code{runs/<case>/phase3/} and \code{runs/<case>/phase4/}, each with a \code{*\_meta.json} recording package versions and the git commit. Trajectory states are in \code{.cache/trajectory/}. Seeds: Step 4 and Phase 2 snapshots 2, Phase 4 7, Phase 3 5.

\end{document}
